% Sources: CHRISTODOULOU_FAITHFUL_PLAN.md (with deviations); CHRISTODOULOU_REPLICATION.md (faithful run);
% prism/analyses/_christodoulou_vendor.py; refs christodoulou2026, vangriethuysen2017.
\paragraph*{The published radiomics pipeline.} We compared PRISM's imaging model with the radiomics pipeline published
on MU-Glioma-Post \citep{christodoulou2026}, run with its authors' public code and settings on the same patients,
masks, label and folds as PRISM. The pipeline describes the tumour region of the current scan by 4,892 radiomic
features computed with PyRadiomics \citep{vangriethuysen2017}: measures of shape, intensity and texture in each of the
four sequences, taken from the original image and from Laplacian of Gaussian and wavelet filtered versions of it.
Within each training set it ranks the features, keeps the top 256 and fits a gradient boosted tree model (LightGBM)
that weights the two classes equally. It tunes the model's settings over 400 trials of a Bayesian search (Optuna),
each trial scored by fivefold cross validation grouped by patient within the training set, and maps the model's
output to a probability by Platt calibration, a logistic fit to the held-out predictions of that cross validation.
The authors' ranking combines the size of L1 penalised logistic regression coefficients with permutation importance,
but their code computes permutation importance only in inner folds of at least 30 scans. Our inner folds held 13 to 16
scans, so the features were ranked by coefficient size alone. We left the code unchanged.
\todo{SUPPLEMENT: extraction settings, the first re-implementation (sensitivity analysis) and the run with the
authors' patient-level label (exploratory).}
The authors' code labels every scan of a patient who ever progressed as progression. We kept PRISM's label, the scan
at which progression was documented, so the comparison is of two methods on one task, not a replication of the
published estimate. The plan, committed before the run, stated that if the pipeline did not rank below PRISM's
imaging model, we would withdraw the claim that the imaging model ranks scans better. We also ran the pipeline on the
masks of a segmentation network, as part of the segmentation comparison.

% Sources: MASK_SENSITIVITY_PLAN.md; METHODS_RESULTS.md E5; CHRISTODOULOU_BRATS_MASKS_PLAN.md (with deviations);
% SEGMENTER_STAGE_PLAN.md; SEGMENTER_STAGE.md; SEGMENTER_BRATS24_PLAN.md.
\paragraph*{Segmentation comparison.} Because the measurements depend on whose outlines are used, we tested how much
PRISM's results change when the masks supplied with each dataset are replaced by a segmentation network's, in three
planned analyses. First, at MU we replaced every supplied mask with the Swin UNETR label map and refitted the imaging,
combined and time-and-enhancement models in the same folds. Swin UNETR does not mark the resection cavity, so the
cavity measurements kept the supplied mask. The published radiomics pipeline was run on the same label maps, with the
authors' tumour region (every labelled voxel) and nothing else changed. Second, we ran PRISM from Swin UNETR alone at
all three centres, taking nothing from the supplied masks. Without the cavity, 22 of the 27 measurements remain. We
refitted the imaging, reduced combined and time-and-enhancement models on MU and applied them without retraining at
LUMIERE and UCSF-ALPTDG, where the time-and-enhancement model used the approximated window. This plan was committed
before any segmenter-derived number was compared with a label at LUMIERE or UCSF-ALPTDG. Third, the BraTS 2024 model
takes the place of Swin UNETR in the second analysis. It marks the cavity, so all 27 measurements come from it.

Each analysis compares every score with its value on the supplied masks by the paired AUC difference. The first
analysis required every paired CI to include zero. The second and third required every paired CI to include zero or
lie above it, so that a segmenter may raise a score but not lower it. All three required the combined model's gain
over the time-and-enhancement model to be read the same way under both mask sets (a gain, or no clear gain), at MU in
the first analysis and at MU and LUMIERE in the others. For the published pipeline, the plan required it to rank below
PRISM's imaging model on the Swin UNETR masks as well; otherwise the comparison would be limited to the supplied
masks. We also report how well each segmenter agrees with the supplied masks, by the Dice overlap (0 for disjoint
outlines, 1 for identical ones), and how noisy the enhancing volume is, as the median fold change between consecutive
scans at most 30 days apart at MU and LUMIERE. Before the run we expected one network applied to every scan to be at
least as repeatable as the supplied masks.

% Sources: NUMBERS.md sections 6, 9, 13, 14; METHODS_PLAN.md; CLINICIAN_BASELINE_PLAN.md; REVIEW_ANALYSES_PLAN.md.
\paragraph*{Tests fixed in advance.} Table~\ref{tab:2} lists each planned test, the rule fixed before it was run and
the outcome, including the two tests that passed only after a revision. The primary results are the external AUCs,
the gain over the time-and-enhancement model and the comparison with the published radiomics pipeline. The other
tests each answer a separate question with its own rule, and we made no correction for multiple testing. Analyses that
were not fixed in advance are marked exploratory where they appear.

% Sources: NUMBERS.md section 9; METHODS_RESULTS.md; CHRISTODOULOU_REPLICATION.md; SEGMENTER_STAGE.md.
\begin{table}[!htbp]
\centering
\caption[Planned tests]{\textbf{Planned tests and their outcome.} Each test, the rule committed to the repository
before it was run, and the outcome. Differences are in AUC, with 95\% CIs.}
\label{tab:2}
\begin{threeparttable}
\footnotesize
\zebrarows
\begin{tabularx}{\linewidth}{@{}R{4.0cm}R{3.9cm}LR{0.9cm}@{}}
\toprule
Test & Rule fixed in advance & Outcome & Section \\
\midrule
Published radiomics pipeline against the clinical-history model on MU folds & Estimate only & $-$0.24 ($-$0.34 to $-$0.15) & \ref{sec:3.1} \\
Published radiomics pipeline below the imaging model on MU folds & Otherwise the comparison is withdrawn & \textbf{Yes}. $-$0.13 ($-$0.23 to $-$0.02) & \ref{sec:3.1} \\
MU results unchanged with Swin UNETR masks (supplied cavity kept) & Every paired CI includes 0, time-and-enhancement reading unchanged & \textbf{No} for the imaging model's AUC. \textbf{Yes} for the combined model, the time-and-enhancement model and the reading & \ref{sec:3.3} \\
Published radiomics pipeline below the imaging model on the Swin UNETR masks & Otherwise the claim is limited to the supplied masks & \textbf{Yes} & \ref{sec:3.1} \\
Results at three centres unchanged with PRISM's Swin UNETR stage & Every paired CI includes 0 or lies above it, gain reading unchanged at MU and LUMIERE & \textbf{No} for the MU imaging model. \textbf{Yes} for the other eight scores and the reading & \ref{sec:3.3} \\
Results at three centres unchanged with the BraTS 2024 stage & As for Swin UNETR & \todo{S4 PENDING} & \ref{sec:3.3} \\
Image measurements and schedule inputs add to the clinical-history model on MU & Gain $\geq$ 0.03, CI above 0 & \textbf{No}. $+$0.02 ($-$0.03 to $+$0.06) & \ref{sec:3.2} \\
MU AUCs unchanged with all scans before progression as negatives & Each changes by $<$ 0.05, time-and-enhancement reading unchanged & \textbf{No} for the imaging model, $-$0.054. \textbf{Yes} for the others and the reading & \ref{sec:3.1} \\
Unknown IDH and MGMT status as its own category & Every paired CI includes 0 & \textbf{Yes} & \ref{sec:2.2} \\
Deep networks tuned on MU add to the clinical-history model with image measurements & Gain $\geq$ 0.03, CI above 0 & \textbf{No}. No network gained & \ref{sec:3.2} \\
Combined model at LUMIERE & AUC $\geq$ 0.65 & \textbf{Yes} & \ref{sec:3.1} \\
Imaging model at UCSF & AUC $\geq$ 0.65, lower bound $>$ 0.5 & \textbf{Yes} & \ref{sec:3.1} \\
Disagreement of the two models marks scans the combined score ranks less well & AUC difference, CI above 0 & \textbf{Yes} at MU; not at LUMIERE or UCSF & \ref{sec:3.4} \\
Combined model over the time-and-enhancement model & CI above 0 & \textbf{No clear gain} at MU or LUMIERE & \ref{sec:3.2} \\
Marginal conformal sets, combined model & Coverage $\geq$ 0.78 in each cohort & \textbf{Yes} on point estimates & \ref{sec:3.4} \\
Reports pass the checker & $\geq$ 95\% in each cohort & \textbf{No} as planned (every LUMIERE report failed). \textbf{Yes} after one revision of the record fields and the prompt & \ref{sec:3.4} \\
Verifier closes a blind spot & $\geq$ 95\% of synthetic errors of that kind rejected & After one revision of the rules, \textbf{Yes} for direction and swapped volumes, \textbf{No} for a swapped tier word & \ref{sec:3.4} \\
Usefulness to clinicians & Not tested & Future work & \ref{sec:4} \\
\bottomrule
\end{tabularx}
\begin{tablenotes}\footnotesize
\item Each plan and the date it was committed are listed in Supplementary Section~S5. The conformal sets of the imaging model at LUMIERE (0.71) and the clinical-history
model at UCSF (0.70) did not meet the coverage rule.
\end{tablenotes}
\end{threeparttable}
\end{table}

% Sources: NUMBERS.md 2.3; SYSTEM_RESULTS.md; METHODS_RESULTS.md (bootstrap seeds, E6a); prism/analyses/detection.py
% (bootstrap, ci); ref vancalster2019.
\paragraph*{Statistics.} Every AUC carries a 95\% confidence interval (CI) from a bootstrap of patients. Each bootstrap
sample draws patients with replacement and keeps all scans of each patient drawn, and the CI is the 2.5th to 97.5th
percentile of the AUC over 1,000 samples. Drawing patients rather than scans keeps a patient's scans together, so the
CI reflects the number of patients, not the larger number of scans. The bootstrap resamples the held-out predictions
without refitting the models, so the CI shows the uncertainty from the patients evaluated, not from model fitting. At
MU each sample averages the AUC over the five repeats of cross validation, as the point estimate does. Analyses added
later used 2,000 samples: the time-and-enhancement model, the AUC split, the group-conditional sets, the segmentation
comparison, and the analyses of calibration, of all scans before progression and of unknown marker status. A
difference between two scores is computed within the same samples, so variation common to both scores cancels. We
read differences from their CIs and report no P values: a difference whose CI lies above zero counts as a gain, and
one whose CI includes zero as no clear difference, not as evidence that the scores are equal. Calibration, the match
between predicted probabilities and observed rates, is summarised by the calibration intercept (above 0 when the
score predicts too low on average) and the calibration slope (below 1 when the predictions are too extreme, above 1
when they are too compressed) \citep{vancalster2019}.
